% Generated draft: requires mathematical and source review.
\documentclass[11pt]{article}
\usepackage[margin=25mm]{geometry}
\usepackage{fontspec}
\setmainfont{Latin Modern Roman}
\usepackage{amsmath,amssymb,mathtools,unicode-math}
\setmathfont{Latin Modern Math}
\usepackage{enumitem,needspace,xurl,hyperref}
\setlength{\parindent}{0pt}
\setlength{\parskip}{5pt}
\setlength{\emergencystretch}{3em}
\newcommand{\N}{\mathbb{N}}
\newcommand{\Z}{\mathbb{Z}}
\newcommand{\Q}{\mathbb{Q}}
\newcommand{\R}{\mathbb{R}}
\newcommand{\C}{\mathbb{C}}
\newcommand{\F}{\mathbb{F}}
\newcommand{\A}{\mathbb{A}}
\newcommand{\PP}{\mathbb{P}}
\newcommand{\E}{\mathbb{E}}
\newcommand{\eps}{\varepsilon}
\newcommand{\dd}{\,\mathrm{d}}
\begin{document}
\section*{UnsolvedMath ID: 10400058}
\textbf{Source problem code:} catalogue entry
\subsection{Definitions and mathematical statement}
Let \(O\) be the trivial knot and \(Z(O)\) its Kontsevich invariant in the source's completed diagram algebra. Its normalization and framing convention are part of the data; one must not automatically replace \(Z(O)\) by \(1\). Let \(\iota\) be the diagram-to-arrow-diagram map referenced in Conjecture 2.16, with the intended extension to this completed setting.

The right-hand side is a closure of a formal exponential with coefficient \(1/2\) and a difference of two pictured diagrams. Label those missing pictures \(D_1,D_2\), solely as placeholders. The recoverable structure of the conjecture is
\[
\iota(Z(O))=
\operatorname{cl}\left(\exp\bigl(\tfrac12(D_1-D_2)\bigr)\right).
\]
In a specified complete characteristic-zero diagram algebra and for a positive-degree argument, the formal exponential means
\[
\exp\bigl(\tfrac12(D_1-D_2)\bigr)
=\sum_{k\ge0}\frac{(D_1-D_2)^k}{2^k k!},
\]
using the source's product. Completion makes this a degreewise formal expression, not a numerical analytic-convergence claim.

The extraction contains neither \(D_1,D_2\) nor the precise map/product/closure operations. The placeholders are not arbitrary diagrams whose choice could be optimized to make the equation true. Closing a skeleton and taking topological closure are also different operations; the source's intended closure must be recovered. Since Conjecture 2.16 was itself pictured, its interval/circle and coefficient conventions cannot be inferred here.

Thus the task is to prove the specific diagram identity of Bar-Natan--Haviv once those missing pictures and conventions are restored, not merely an ordinary formula for the unknot's Kontsevich invariant. The supplied research summary identifies the conjecture but supplies no resolved-status theorem or complete formula.
\subsection{Short English statement}
Prove the stated identity for the arrow-diagram image of the unknot's Kontsevich invariant. Its two defining diagrams, closure/product operations, map extension, and normalization are missing from the extracted formula.
\subsection{Sources}
\begin{itemize}
\item Local UnsolvedMath registry, permanent ID \(10400058\), exact source code \texttt{catalogue entry}, original Conjecture 3.8, attributed to D. Bar-Natan and A. Haviv.
\item Ohtsuki, \emph{Problems on invariants of knots and 3-manifolds} (2002), Conjecture 3.8, recorded PDF page 55; \url{https://msp.org/gtm/2002/04/gtm-2002-04-024s.pdf}. Recorded access 2026-07-29; pdf-font-structure extraction.
\item Source cross-reference Conjecture 2.16 for \(\iota\). The Kontsevich-invariant citation is extracted as ``[3 5]''; whether this is [35] or separate references cannot be recovered from the record.
\end{itemize}
\end{document}
